\documentclass[a4paper, UTF-8]{article}
\usepackage{algorithm}
\usepackage{algpseudocode}
\usepackage{amsmath}
\usepackage{amsfonts}
\usepackage{amsthm}
\usepackage{tikz}
\usetikzlibrary{automata, positioning, arrows, trees}
\usepackage[utf8]{inputenc}
\usepackage[T1]{fontenc}
\usepackage{hyperref}
\usepackage{url}
\usepackage{booktabs}
\usepackage{nicefrac}
\usepackage{microtype}
\usepackage{xcolor}
\usepackage{amssymb}

\theoremstyle{plain}
\newtheorem{theorem}{\indent Theorem}[section]
\newtheorem{lemma}[theorem]{\indent Lemma}
\newtheorem{assumption}{\indent Assumption}[section]
\newtheorem{note}[theorem]{\indent Notation}
\newtheorem{proposition}[theorem]{\indent Proposition}
\newtheorem{corollary}[theorem]{\indent Corollary}
\newtheorem{definition}{\indent Definition}[section]
\newtheorem{example}{\indent Example}[section]
\newtheorem{remark}{\indent Remark}[section]
\newenvironment{solution}{\begin{proof}[\indent\bf Solution]}{\end{proof}}
\renewcommand{\proofname}{\indent\bf Proof}

\begin{document}

\title{\bf{Lab 10 Answers}}
\author{Yanqiao Chen*\\ 
  Department of Computer Science and Engineering\\
  Southern University of Science and Technology\\
  \texttt{12412115@mail.sustech.edu.cn}
  }
\maketitle

\section{Question 1} 

We show that $\overline{EQ_{CFG}}$ is Turing-recognizable. We construct the following recognizer. 

$$
D = \text{``On input } \langle G_1, G_2 \rangle\text{, where } G_1 \text{ and } G_2 \text{ are CFGs:}
$$

\begin{itemize} 
  \item Enumerate all strings in $\Sigma^*$ in lexicographic order, and for each string $w$, do the following:
  \begin{itemize}
    \item Run the Turing machine for $A_{CFG}$ on input $\langle G_1, w \rangle$ and the Turing machine for $A_{CFG}$ on input $\langle G_2, w \rangle$ in parallel. 
    \item If one of the two Turing machines accepts and the other rejects, then accept.
  \end{itemize} 
\end{itemize}

If they are different, then there must exist a string $w$ that is generated by one of the two CFGs but not the other. In this case, $D$ will accept. If they are the same, then they may loop or reject, but they will never accept. Therefore, $D$ is a recognizer for $\overline{EQ_{CFG}}$.

Thus, $\overline{EQ_{CFG}}$ is Turing-recognizable, thus $EQ_{CFG}$ is co-Turing-recognizable.

\section{Question 2} 

No, it does not. Let $f(x) = \begin{cases} 1, & x\in A \\ 0, & x\notin A\end{cases}$. Thus let $B = \{1\}$. Then $f$ is a reduction from $A$ to $B$. However, let $A = \{1^k0^k \mid k \geq 0\}$, which is not regular, and $B = \{1\}$, which is regular. Therefore, $A$ is not regular, but $B$ is regular. Thus, the regularity of $B$ does not imply the regularity of $A$.

\section{Question 3} 

We know $A$ is decidable $\iff$ $A$ is Turing-recognizable and $\overline{A}$ is Turing-recognizable. Since $A$ is Turing-recognizable, we only need to show that $\overline{A}$ is Turing-recognizable and $A \leq_{m} \overline{A}$. Thus, $f(x) \in \overline{A}$ if and only if $x \in A$ and $f(x) \in A$ if and only if $x \in \overline{A}$. We construct the following Turing machine to recognize $\overline{A}$.

$$
D = \text{``On input } w\text{, where } w \text{ is a string:}
$$

\begin{itemize}
  \item For any input, $D$ run the Turing machine to map $w$ to $f(w)$, and then run the Turing machine for $A$ on input $f(w)$.
  \item If the Turing machine for $A$ accepts, then reject; if it rejects, then accept.
\end{itemize}

Thus, $\overline{A}$ is Turing-recognizable. Therefore, $A$ is decidable.

\section{Question 4} 

Let $M'$ be a Turing machine. 
$$
M' = \text{``On input } \langle x \rangle\text{, where } x \text{ is a string:}
$$

\begin{itemize}
  \item If $x = 01$, accept.
  \item Then run $M$ on input $w$, if $M$ accepts, then accept; if $M$ rejects, then reject.
\end{itemize}

This construction perform a mapping reduction $f(\langle M, x \rangle) = \langle M' \rangle$. If the decider decides $M'$ whether is accept $w$ if and only if it accepts $w^{R}$, then if $x \ne 01$, then $M'$ accepts $x$ if and only if $M$ accepts $w$, which means that it accept both $w$ and $w^{R}$. Thus, if a decider can decide $M' \in T$, then it must decide whether $M'$ accept $w^{R}$, which include $10$, thus it must decide whether $M$ accepts $w$, which is equivalent to deciding $A_{TM}$.  

Thus, if $M'$ is decidable, then the decider decides $A_{TM}$, which is a contradiction. Therefore, $T$ is undecidable.

\section{Question 5} 

We prove from right to left. If $A \leq_m A_{TM}$, we construct the following Turing machine to decide $A$.

$$
D = \text{``On input } w\text{, where } w \text{ is a string:}
$$

\begin{itemize}
  \item Run the Turing machine to map $w$ to $f(w)$, and then run the Turing machine $M$ on input $\langle M, f(w) \rangle$. Then run the Turing machine for $A_{TM}$ on input $\langle M, f(w) \rangle$ to recognize whether $M$ accepts $f(w)$.
  \item If the Turing machine for $A_{TM}$ accepts, then accept; if it rejects, then reject.
\end{itemize} 

We prove from left to right. If $A$ is decidable, we construct the following Turing machine to perform a mapping reduction from $A$ to $A_{TM}$. As $A$ is T-recognizable, there exists a Turing machine $M$ that recognizes $A$. Let $f(x) = \langle M, x \rangle$, then we construct the following Turing machine to perform a mapping reduction from $A$ to $A_{TM}$: 
$$
D = \text{``On input } \langle M, w\rangle \text{, where } w \text{ is a string, }M \text{ is a Turing machine:}
$$

\begin{itemize}
  \item Simulate the Turing machine for $A$ on input $w$. If it accepts, then accept 
  \item If it rejects, then reject.
\end{itemize}

Thus $A \leq_m A_{TM}$, which completes the proof.

\section{Question 6} 

Construct a mapping reduction from $A_{TM}$ to T. 
$$
f(\langle M, w \rangle) = \langle M' \rangle
$$ 

where $M'$ is a Turing machine defined as follows:
$$
M' = \text{``On input } x \text{, where } x \text{ is a string:}
$$
\begin{itemize}
  \item If $x \ne \epsilon$, accept.
  \item Then run $M$ on input $w$, if $M$ halts on $w$, accept; otherwise, reject.
\end{itemize}

Thus, if $M'$ can be decided, then we can decide whether $M$ accepts $w$, which is equivalent to deciding $HALT_{TM}$. Therefore, $T$ is undecidable.

\section{Question 7} 

\begin{itemize}
  \item Undecidable. "$L(M)$ is infinite" is non-trivial, thus it is undecidable by Rice's theorem. 
  \item Undecidable. "$1011 \in L(M)$" is non-trivial, thus it is undecidable by Rice's theorem.
  \item Undecidable. "$L(M) = \Sigma^*$" is non-trivial, thus it is undecidable by Rice's theorem.
\end{itemize} 

\section{Question 8} 

\begin{itemize}
  \item True. 
  \item False. 
  \item True. 
  \item False. 
  \item False. 
  \item True. 
\end{itemize}

\section{Question 9}




\end{document}